\section{Design}
\label{sec:design}

\subsection{Logical Topology Construction}

The software back-end starts from a topology specification that lists enabled
logical endpoints, their virtual BDFs, device templates, and backend bindings.
It constructs Type-1 bridge images and Type-0 endpoint images, including vendor
and device identity, class code, BAR shape, and the capabilities that \sys can
faithfully implement. Capabilities are promises: for example, the current NIC
template omits MSI-X because the prototype implements a shared level interrupt
rather than a virtual MSI-X table and vector remapper.

Topology realization produces two coupled artifacts. A compact ECAM image is
copied into FPGA-visible shadow memory, and a route table maps assigned BAR
ranges to logical backend identifiers. A readiness protocol hides both paths
until the image and routes are initialized, preventing Linux enumeration from
observing a partially constructed hierarchy.

\subsection{Configuration-Space Virtualization}

The FPGA front-end exposes a standard ECAM aperture. Reads to defined BDFs use
the shadow memory fast path; reads to absent functions return all ones. Writes
become fixed-size control events carrying BDF, offset, byte enables, payload,
and a sequence number. Software applies PCI configuration semantics to the
authoritative image, recomputes any BAR route affected by probing or assignment,
writes the changed state back to the shadow, and acknowledges the matching
sequence. Only then does the FPGA complete the DUT write. This rule makes a
subsequent read and BAR access observe the new state.

The shadow is therefore a read cache rather than an independent device model.
Software owns writable semantics, while the FPGA supplies predictable ECAM
latency and a standard interface to the DUT.

\subsection{Endpoint Semantic Mediation}

For runtime access, the FPGA matches a BAR address against its route entries and
emits an event tagged with logical BDF, BAR, offset, width, byte enables, and
sequence. The generic manager dispatches the event to a per-type semantic
adapter. Reads return a value and status through a mailbox. Ordinary writes can
be acknowledged after the required backend effect. Submission writes require a
stronger ordering rule.

A native driver prepares commands or descriptors in memory before writing a
doorbell. The physical backend cannot consume the original entry because it
contains DUT addresses. The adapter waits until the FPGA has accepted the BAR
write and the associated memory entry is stable, parses the entry, translates
its DMA references, writes the translated form, and only then forwards the
physical doorbell. This sequence preserves ownership transfer without requiring
the DUT driver to understand the proxy.

Device semantics remain explicit. The NVMe adapter understands controller
registers, queue setup, doorbells, PRPs and PRP lists, and phase-tagged CQEs. The
NIC adapter understands ring bases and lengths, descriptor buffer addresses,
tail updates, reset behavior, and interrupt cause/mask registers. Common
dispatch and transport do not make unlike devices semantically interchangeable.

\subsection{DMA Data-Path Realization}

\sys separates address retargeting from byte movement. On submission, a semantic
adapter replaces each DUT physical address with an address reachable by the
selected physical backend. For host-local devices, the translated address can
target a coherent alias of DUT memory through the FPGA DMA aperture, allowing
the physical device to move payload without a software copy. A remote backend
may instead use staging or a transport-specific direct path. Both modes preserve
the same logical endpoint contract even though their performance and copy count
differ.

Completion follows the reverse dependency. The back-end first makes returned
payload and device writebacks visible in DUT memory. It then publishes the
guest-visible completion state, updates the adapter's pending state, and finally
changes the virtual interrupt level. This ordering prevents a driver from being
woken before the data it will consume.

\subsection{Interrupt Virtualization}

The initial prototype exports legacy INTx. Each backend maintains private
pending state, and the manager ORs those states into one aggregate level. A
queue-consumption doorbell, mask update, or reset changes only the owning
backend; the physical line is deasserted only when no backend remains pending.
For the NIC, the adapter also preserves read-clear and set/clear semantics of
the interrupt cause and mask registers. MSI/MSI-X requires per-vector table,
address, and delivery virtualization and is outside the current implementation.

\subsection{Scaling the FPGA Mechanism}

The FPGA contains endpoint-independent mechanisms: one ECAM service path, one
control-event channel, one mailbox, a parameterized route table, DMA windows,
and interrupt injection. Adding another supported endpoint consumes another
configuration slot, route state, and software object, but does not replicate an
NVMe or NIC controller in RTL. Adding a device type primarily extends the
software semantic adapter. The remaining hardware cost is real and must be
measured; the intended hypothesis is that it grows with generic capacity rather
than with duplicated device protocol engines. The 1/2/4/8/13-slot synthesis
sweep in Section~\ref{sec:evaluation} is the experiment that can turn this
hypothesis into a quantitative scalability claim.
